\section{Implementation and Experimental Evaluation}
\label{sec:evaluation}

This section details the implementation of the MPTA-Repair framework
and systematically evaluates its practical effectiveness across both
standardized benchmark models and complex industrial case studies. To
clearly distinguish our targeted contribution while maintaining
practical applicability, our methodology strictly focuses on the
automated repair of numerical clock bounds within transition guards and
location invariants, deliberately avoiding arbitrary alterations to the
underlying discrete transition logic. This focused parametric repair
space is pragmatic for industrial contexts, because modifying clock
bounds can correct clerical modeling mistakes and provide useful
insights for redimensioning system time resources without disrupting the
intended functional design. Furthermore, we compare MPTA-Repair with a
specialized iterative baseline technique constructed by extending the
foundational TARTAR repair idea \cite{koelbl2020tartar} with a
sequential repair loop. This comparison demonstrates the necessity of
incorporating multi-property regression and multi-counterexample
analysis to prevent secondary errors during repair.

\subsection{Prototype Implementation}

We implemented MPTA-Repair as a highly integrated, incremental repair
pipeline built around a specialized timed-automata model parser, a
verification interface using UPPAAL \cite{larsen1997uppaal}, a
MaxSMT-based repair backend powered by Z3 \cite{demoura2008z3}, a
dedicated admissibility checker, and an overarching iterative refinement
control loop. Initially, the repair frontend parses the faulty model to
extract editable clock-bound variables embedded within transition guards
and location invariants. The trace analyzer then maps observed
diagnostic traces to these extracted variables. Concurrently, the
relation analyzer establishes dependencies among the specified
properties, generated traces, and localized repair sites in order to
prioritize candidate generation and reduce the computational search
space.

Once the MaxSMT backend synthesizes a viable candidate assignment, the
model writer instantiates the corresponding repaired model. Final
candidate acceptance relies on two consecutive validation gates designed
to guarantee both timing correctness and behavioral fidelity. The first
gate performs full-property regression verification through UPPAAL to
ensure that no collateral property violations occur. The second gate is
an admissibility check for functional equivalence following the
TARTAR-style criterion \cite{koelbl2019clock}. This module constructs
untimed automata from both the original and repaired models and performs
a comparative language-equivalence test. The system rejects any
candidate that inadvertently adds or removes functional untimed
behavior. Whenever either validation gate rejects a candidate, the
control loop captures the newly exposed diagnostic traces and feeds them
back into the relation analyzer. This initiates an incremental,
counterexample-guided refinement iteration that further constrains the
MaxSMT solver with updated evidence until a globally correct and
admissible repair is found or the search budget is exhausted.

\subsection{Evaluation Strategy and Datasets}

We evaluate the proposed framework on two dataset groups. The benchmark
dataset consists of timed-automata models collected from UPPAAL and XTA
benchmark sources. These models provide structural diversity and allow
comparison with existing clock-bound repair settings. The industry
dataset contains practical case studies involving railway alarm and
communication scenarios
\cite{gonzalez2024realtimedevs,basile2021unisig}, train-gate behavior
\cite{behrmann2004tutorial}, automotive gear control
\cite{lindahl1998gear}, cardiac pacing \cite{jiang2012pacemaker},
autonomous road-junction rules \cite{belguidoum2021junction},
mechanical ventilation \cite{cuartas2023ventilator}, and aerospace
scheduling \cite{david2010herschel}. Together, the two dataset groups
test both algorithmic behavior on established examples and practical
behavior on industrial models with synchronization, multiple
components, and domain-level timing requirements.

Because naturally occurring faulty timed-automata models are scarce, we
use a mutation-based evaluation strategy \cite{jia2011mutation}. Faults
are seeded into verified reference models by modifying clock-bound
constants in guards and invariants. For simple mutants, one bound is
changed and the resulting model is retained if it violates at least one
selected timing safety property. For complex mutants, multiple
clock-bound edits are injected to simulate interacting timing faults.
Generated mutants are filtered to remove invalid or trivial cases: each
retained mutant must be syntactically valid and must violate at least
one selected property.

We compare MPTA-Repair with an adapted iterative TARTAR-style baseline
\cite{koelbl2020tartar}. For fairness, the baseline receives the same
full-property regression and admissibility validation gates used by
MPTA-Repair. However, the baseline generates repairs from one violated
property or diagnostic trace at a time. It does not jointly analyze
relations among multiple properties, counterexamples, and repair
variables before candidate generation. This difference allows us to
measure the benefit of the multi-property and multi-counterexample
repair strategy.

\subsection{Experimental Results}

The experimental results show that MPTA-Repair outperforms the
iterative baseline, particularly on complex industrial models. On the
benchmark dataset, the baseline repairs 72\% of the cases, while
MPTA-Repair repairs 93\%. On the industry dataset, the gap is larger:
the baseline repairs 20\% of the cases, while MPTA-Repair repairs
76\%. These results suggest that the proposed workflow remains
effective on established timed-automata benchmarks and is more robust
when properties and traces interact in industrial models.

\begin{table}[t]
\caption{Overall repair outcomes reported in the source draft.}
\label{tab:overall-results}
\centering
\scriptsize
\setlength{\tabcolsep}{2pt}
\begin{tabular}{@{}p{0.24\columnwidth}p{0.34\columnwidth}p{0.09\columnwidth}p{0.11\columnwidth}p{0.13\columnwidth}@{}}
\toprule
Dataset group & Method & Cases & Repaired & Success rate \\
\midrule
Benchmark & Iterative TARTAR-style baseline & N1 & R1 & 72\% \\
Benchmark & MPTA-Repair & N1 & R2 & 93\% \\
Industry & Iterative TARTAR-style baseline & N2 & R3 & 20\% \\
Industry & MPTA-Repair & N2 & R4 & 76\% \\
\bottomrule
\end{tabular}
% TODO: Replace placeholder counts N1, N2, R1, R2, R3, and R4 with final evaluated counts, and confirm the success rates.
\end{table}


The baseline remains competitive on simpler benchmark models where
faults are highly localized. In such cases, a single diagnostic trace
often provides enough evidence to infer a useful clock-bound edit.
However, in industrial models with interacting concurrent components,
synchronization mechanisms, and observer templates, local baseline
repairs frequently fail later global validation checks.

Our analysis identifies two main reasons for the stronger performance
of MPTA-Repair. First, multi-counterexample accumulation reduces
overfitting to one timed diagnostic trace. A single trace usually
exposes only one timing realization of a deeper fault, so a candidate
that eliminates that trace may still leave related violations feasible.
By accumulating diagnostic traces from multiple operational modes,
MPTA-Repair can synthesize small joint edit sets that address the shared
timing fault without requiring excessive refinement rounds.

Second, relation-guided optimization improves search efficiency and
candidate quality. Property relations help avoid repeatedly solving
equivalent obligations. Counterexample relations identify traces that
share path fragments or repair variables. Repair-variable relations
focus the solver on clock bounds close to the observed violations
rather than opening all syntactically editable bounds in the model. In
the ablation analysis, removing relation guidance increased the number
of generated candidates and the frequency of regression failures.
Successful repairs typically modified only one or two clock bounds, and
the average edit magnitude remained close to the injected fault size,
which makes the suggested repairs easier for engineers to inspect.

\subsection{Discussion and Practical Lessons}

The evaluation highlights several practical lessons for applying
automated repair to industrial timed automata. First, full-property
regression testing is necessary. Industrial properties are often
interdependent: a local fix that satisfies a response monitor can
violate a separate safety monitor that constrains maximum waiting time.
MPTA-Repair avoids accepting such regressions by treating restoration of
the full property set as a global objective rather than as a sequence of
isolated local patches.

Second, admissibility changes the meaning of repair success. Without
checking functional equivalence, a solver could satisfy a timed property
by disabling intended behavior, for example by blocking a necessary
communication response or preventing a critical action. Such a candidate
may eliminate the observed timed diagnostic trace, but it is not a valid
repair for an industrial model. The admissibility check therefore
prevents property satisfaction from being achieved through unacceptable
loss of untimed behavior.

Finally, the unrepaired cases show the limits of the current repair
space. Some attempts fail because of solver, verification, or
admissibility timeouts, and some models use syntax outside the supported
fragment. Other cases appear to lack a feasible repair within numerical
clock-bound edits alone. These failures suggest that broader repair
actions, such as changing synchronization labels, clock resets, or
discrete data updates, are an important direction for future work.
